\section{Thesis Outline}
\label{sec:thesis-outline}

The remainder of this thesis is organised as follows.

\textbf{\Cref{ch:key-concepts-related-work}}, \emph{Key Concepts and Related Work}, surveys the concepts and literature that underpin the work.
It defines the key terms used throughout the thesis, reviews related work on eye tracking in reading, adaptive reading systems, and the existing market landscape, and identifies the gaps that this thesis addresses.
The chapter closes with the refined problem statement and four research questions that govern the rest of the work.

\textbf{\Cref{ch:methodology}}, \emph{Methodology}, describes how the thesis was conducted.
It positions the work within the Design Science Research paradigm, explains how requirements were elicited, describes the iterative development process, sets out the evaluation strategy, and addresses ethical considerations arising from the use of eye-tracking data.

\textbf{\Cref{ch:requirements}}, \emph{Requirements and Analysis}, derives the functional and non-functional requirements of the platform from stakeholder analysis and use cases.
It introduces a domain model and documents the constraints that bound the design space.

\textbf{\Cref{ch:design}}, \emph{System Design and Architecture}, presents the system architecture.
It introduces the four-layer modular design, describes the module contracts and data flow, and records the key design decisions and their rationale.

\textbf{\Cref{ch:implementation}}, \emph{Implementation}, describes how the architecture is realised in code.
It covers the monorepo structure, selected implementation highlights across the sensing, decision, and intervention layers, and the build and CI workflow.

\textbf{\Cref{ch:evaluation}}, \emph{Evaluation}, evaluates the platform against the four research questions.
It reports on capability demonstration, calibration validation, sensing pipeline latency, and session data export analysis, and assesses the degree to which the platform meets each criterion.

\textbf{\Cref{ch:discussion}}, \emph{Discussion}, interprets the evaluation results in relation to the problem statement and the literature, discusses the limitations of the current work, and outlines directions for future research.

\textbf{\Cref{ch:conclusion}}, \emph{Conclusion}, summarises the contributions and answers the research questions directly.
It introduces no new content.

Two appendices document the artifact for a reader who would run or operate it.
App.~\ref{app:installation-guide} is an installation guide that covers setup from a clean checkout to a verified system, and App.~\ref{app:user-guide} is a user guide for the researcher operating a reading session.
